\documentclass[twocolumn]{article}
\usepackage[utf8]{inputenc}

\counterwithout{figure}{section} % Si usas la clase 'article'
\usepackage{graphicx} % graficos
\usepackage{amsmath} %matrices
\usepackage{float}    
\usepackage{hyperref}


\graphicspath{{Imagenes/}} %graficos

\usepackage[spanish]{babel} %idioma

\usepackage{chngcntr}

\spanishdecimal{.}
\setlength{\parskip}{1em} 

\begin{document}
\author{Jesús Armando Cañas Gambo y Miguel Angel Pérdomo Gaitan}
\date{\today}
\title{Efecto fotoélectrico}
\maketitle
\section{Marco Teórico}

\subsection{Efecto fotoélectrico}

El efecto fotoeléctrico consiste en la emisión de eléctrones por un material cuando se lo ilumina con una corriente de fotones.

\end{document}